\section{Conclusion}
\label{sec:conclusion}

We presented Failure-Trace Allocator (FTA) as a deterministic selective-deferral
policy for matched-budget LLM answer selection.
On a disjoint 720-example aggregate evaluation (three independent seeds, zero overlap),
Failure-Trace Allocator (FTA) achieves 80.69\% (581/720) versus L1 77.64\%, S1 77.08\%, and TALE 75.14\%.
Source-stratified 95\% bootstrap confidence intervals are strictly positive for all four
baseline comparisons, with the lower bound vs.\ best external at $+1.11$ pp.
On the Final-300 evaluation (seed 71, budget 6), Failure-Trace Allocator (FTA) reaches 86.67\% (260/300),
compared to 83.00\% for L1, 82.00\% for S1, and 78.33\% for TALE.
Against the pooled four-answer ensemble, Failure-Trace Allocator remains ahead by point estimate but is not statistically separated, so the relevant conclusion is competitiveness with strong aggregation baselines rather than broad superiority over ensembling.

The broader contribution is the framing of post-generation answer selection as a
matched-budget selective-deferral problem over answer sources. FTA is one
deterministic instantiation of that framework, consisting of two rule-based gates
applied in strict priority order.
The Low-Depth Guard (LDG) detects structurally weak frontier search branches using the
logged \texttt{override\_reason} field and substitutes the external majority vote.
The External-Consensus Fallback (ECO) identifies cases where all three external
baselines unanimously disagree with the frontier's selection, overriding to the
consensus answer.
Both gates are gold-free: no correctness labels, example identifiers, or artifact
paths are accessed at runtime.
The rules, trigger conditions, and exact field names are fully specified in
Table~\ref{tab:policy_rules} and Table~\ref{alg:fix24}, with reproducibility
details documented in the appendix and submission package.

The evidence remains bounded to one tested provider/benchmark/budget setting, so
the paper should be read as an answer-selection and selective-routing study rather
than as a claim of a generally superior end-to-end solver. Further accuracy gains
in this setting may require advances in candidate generation rather than selection
policy.
Failure analysis of the 40 residual errors on Final-300 shows that 24 (60\%) are
\textit{all-methods-wrong} cases in which no evaluated method produced the correct answer;
the dominant subtype is multi-step composition traps (10/24) where the error is
embedded in intermediate reasoning rather than in the final answer selection step.
Addressing this residual category may require either stronger base generators or
training-time interventions that improve coverage of hard compositional patterns.
